\HandoutSetup
  {NE 630}
  {02}
  {Nuclear Fission and Chain Reactions}
  {Read FNRP Sections 1.4--1.5. Revisit binding energy per nucleon from Lesson 1 and define the key terms before class.}

\begin{document}
\MakeHandoutTitle

\begin{handoutcolumns}

\begin{fixedbox}{Macro objective}
Characterize the neutron population and energy
produced as functions of time using the neutron lifetime, the energy released
per fission, and the multiplication factor $k$.
\end{fixedbox}

\TermStrip{fusion; fission; fission product; fission-product yield; prompt; delayed; decay heat; fissionable; fissile; fertile; generation; neutron lifetime; multiplication factor; criticality}

\HandoutSection{1}{Stability, fusion, and fission}

\begin{fixedbox}{Binding-energy relations}
Using \emph{nuclear} masses,
\[
 B=\bigl[Zm_p+Nm_n-m_X\bigr]c^2.
\]
Using the \emph{atomic} mass $m({}^{A}_{Z}X)$ and neglecting electron
binding energies,
\[
 B=\bigl[Z\,m({}^{1}_{1}\mathrm H)+Nm_n-m({}^{A}_{Z}X)\bigr]c^2,
 \qquad \frac{B}{A}=\frac{B}{N+Z}.
\]
\end{fixedbox}

\begin{boardbox}{Read the binding-energy curve}
\centering
\begin{tikzpicture}[x=0.0305\linewidth,y=0.23cm,>=Latex]
  \draw[->,KSUGray] (0,0) -- (29,0) node[right,black] {$A$};
  \draw[->,KSUGray] (0,0) -- (0,9.8) node[above,black] {$B/A$};
  
\end{tikzpicture}

\RaggedRight
Why can either process release energy? What competing effects produce the
maximum near the middle of the chart?
\RuledLines{2}
\end{boardbox}


\begin{fixedbox}{Two fusion examples from the notes}
\begin{align*}
{}^{2}_{1}\mathrm H+{}^{3}_{1}\mathrm H
 &\longrightarrow{}^{4}_{2}\mathrm{He}+{}^{1}_{0}\mathrm n+17.59\unit{MeV},\\
{}^{2}_{1}\mathrm H+{}^{2}_{1}\mathrm H
 &\longrightarrow{}^{3}_{1}\mathrm H+{}^{1}_{1}\mathrm H+4.02\unit{MeV}.
\end{align*}
The energy is favorable; bringing two positive nuclei close enough to react
is the engineering challenge.  
For D--T fusion, estimate the Coulomb energy 
$V_C=(4\pi\epsilon_0)^{-1}e^2/r$ at $r\approx 4\unit{fm}$ and convert it to an 
equivalent temperature $T=V_C/k_B$; what does the 
resulting $\approx 0.36\unit{MeV}$ ($\approx 4\times10^9\unit{K}$) tell you about the need for quantum tunneling?
\end{fixedbox}

\begin{fixedbox}{Three forms of nuclear power}
\scriptsize
\textbf{1.}\ \RevealBlank[1.28in]{}\hfill
\textbf{2.}\ \RevealBlank[1.28in]{}\par
\textbf{3.}\ \RevealBlank[3.00in]{}
\normalsize
\end{fixedbox}

\switchcolumn

\HandoutSection{2}{Anatomy of neutron-induced fission}

\begin{fixedbox}{Prompt sequence and time scales}
\centering
\begin{tikzpicture}[
  node distance=0.24in and 0.26in,
  >=Latex,
  every node/.style={font=\scriptsize,align=center}
]

  \node (n) {$n$};

  \node[
    draw,
    circle,
    minimum size=0.38in,
    right=of n
  ] (u5) {$^{235}\mathrm{U}$};

  \node[
    draw,
    circle,
    minimum size=0.43in,
    right=of u5
  ] (u6) {$^{236}\mathrm{U}^*$};

  \node[
    draw,
    ellipse,
    minimum width=0.56in,
    minimum height=0.32in,
    right=of u6
  ] (def) {deform};

  \node[
    minimum width=0.66in,
    minimum height=0.34in,
    right=of def
  ] (neck) {};

  \begin{scope}[shift={(neck.center)}]
    \draw[smooth cycle,tension=0.75]
      plot coordinates {
        (-0.31in,  0.00in)
        (-0.25in,  0.13in)
        (-0.12in,  0.15in)
        (-0.045in, 0.055in)
        ( 0.045in, 0.055in)
        ( 0.12in,  0.15in)
        ( 0.25in,  0.13in)
        ( 0.31in,  0.00in)
        ( 0.25in, -0.13in)
        ( 0.12in, -0.15in)
        ( 0.045in,-0.055in)
        (-0.045in,-0.055in)
        (-0.12in, -0.15in)
        (-0.25in, -0.13in)
      };
  \end{scope}


  \draw[->] (n)    -- (u5);
  \draw[->] (u5)   -- (u6);
  \draw[->] (u6)   -- (def);
  \draw[->] (def)  -- (neck);

\end{tikzpicture}

\RaggedRight
A heavy nucleus usually splits into two energetic fragments. Excitation
energy produces prompt $\gamma$ rays and approximately
\RevealBlank[0.46in]{} neutrons (about
\RevealBlank[0.46in]{} on average in the board sketch).

\smallskip
The split occurs within roughly \RevealBlank[0.66in]{}; prompt
emissions are complete within roughly \RevealBlank[0.66in]{}
of the initiating neutron absorption.
\end{fixedbox}

\begin{boardbox}{Why are fission products neutron rich?}
\begin{minipage}[t]{0.48\linewidth}
\centering
\begin{tikzpicture}[x=0.032\linewidth,y=0.032\linewidth,>=Latex]
  \draw[->,KSUGray] (0,0)--(12,0) node[right,black] {$Z$};
  \draw[->,KSUGray] (0,0)--(0,15) node[above,black] {$N$};
  
\end{tikzpicture}
\end{minipage}\hfill
\begin{minipage}[t]{0.48\linewidth}
\centering
\begin{tikzpicture}[x=0.031\linewidth,y=0.034\linewidth]
  \draw[->,KSUGray] (0,0)--(25,0) node[right,black] {$A$};
  \draw[->,KSUGray] (0,0)--(0,15) node[above,black,rotate=90] {yield};
  
\end{tikzpicture}
\end{minipage}

The parent has a large neutron-to-proton ratio. Splitting it near the usual
two fission-yield peaks leaves fragments on which side of the stability line?
How do they move toward stability?
\RuledLines{2}
\end{boardbox}


\begin{checkpoint}[Prompt versus delayed]
Sort the emissions or energy carriers discussed on the board:
\begin{tabularx}{\linewidth}{@{}>{\bfseries}l X@{}}
prompt & fragment K.E.; \Blank[1.95in]\\[2pt]
delayed & $\beta$; $\gamma$; delayed $n$; \Blank[1.05in]
\end{tabularx}
Which contribution is effectively lost for heat production?
\RevealBlank[1.40in]{}
\end{checkpoint}

\begin{boardbox}{Reading completion: fuel vocabulary}
\scriptsize
\textbf{Fissionable:} \Blank[2.52in]\\[-1pt]
\emph{Examples:} \Blank[2.82in]\\[2pt]
\textbf{Fissile:} \Blank[2.80in]\\[-1pt]
\emph{Examples:} \Blank[2.82in]\\[2pt]
\textbf{Fertile:} \Blank[2.80in]\\[-1pt]
\emph{Pathways:} \Blank[2.72in]
\normalsize
\end{boardbox}

\end{handoutcolumns}

\newpage
\begin{handoutcolumns}

\HandoutSection{3}{Energy release and decay heat}

\begin{boardbox}{Symmetric-fission bookkeeping}
Consider
\[
 n+{}^{235}_{92}\mathrm U
 \longrightarrow 2\,{}^{117}_{46}\mathrm{Pd}+2n.
\]
\scriptsize
\begin{tabularx}{\linewidth}{@{}>{\bfseries}l >{\raggedleft\arraybackslash}X@{}}
$m({}^{235}\mathrm U)$ & $235.043928\unit{u}$\\
$m({}^{117}\mathrm{Pd})$ & $116.917956\unit{u}$\\
$m_n$ & $1.0086649\unit{u}$
\end{tabularx}
\normalsize
{\small
\[
\begin{aligned}
Q_{\mathrm{prompt}}
 &=\bigl[m({}^{235}\mathrm U)-2m({}^{117}\mathrm{Pd})-m_n\bigr]c^2,\\
 &=\RevealBlank[0.92in]{}.
\end{aligned}
\]
}
Show the mass cancellation and arithmetic here:
\RuledLines{2}
\end{boardbox}

\begin{boardbox}{After the neutron-rich fragments decay}
For each fragment,
\[
{}^{117}_{46}\mathrm{Pd}\to{}^{117}_{47}\mathrm{Ag}
\to{}^{117}_{48}\mathrm{Cd}\to{}^{117}_{49}\mathrm{In}
\to{}^{117}_{50}\mathrm{Sn}.
\]
With $m({}^{117}\mathrm{Sn})=116.902955\unit{u}$,
{\small
\[
\begin{aligned}
Q_{\mathrm{total}}
 &=\bigl[m({}^{235}\mathrm U)-2m({}^{117}\mathrm{Sn})-m_n\bigr]c^2,\\
 &=\RevealBlank[0.92in]{}.
\end{aligned}
\]
}
Why can atomic masses be used directly for the complete $\beta^-$ chain?
\RuledLines{1}
\end{boardbox}


\begin{takeawaybox}[Fission-energy scale]
The board estimate is roughly \RevealBlank[0.66in]{} released per
thermal fission of ${}^{235}$U, with about
\RevealBlank[0.72in]{} recoverable. Most prompt energy appears as
fragment kinetic energy; each prompt neutron carries roughly
\RevealBlank[0.56in]{}.
\end{takeawaybox}

\begin{checkpoint}[Decay heat immediately after shutdown]
Delayed $\beta$ and $\gamma$ energy continues as fission products and their
daughters decay. The board note marks an immediate decay-heat level of about
\RevealBlank[0.54in]{} of the pre-shutdown power:
\[
 P_{\mathrm{decay}}(0^+)\approx\MathBlank[0.62in]{}\,P_0.
\]
For a reactor operating at $P_0$, what heat-removal obligation remains after
fission power is terminated?
\RuledLines{1}
\end{checkpoint}

\switchcolumn

\HandoutSection{4}{Multiplication and chain reactions}

\begin{fixedbox}{Generation model}
\[
 k=\frac{n_i}{n_{i-1}},
 \qquad n_i=n_0k^i.
\]
If one generation lasts $\ell$ seconds, then $i=t/\ell$ and
\[
\begin{aligned}
 n(t)&=n_0k^{t/\ell}
     =n_0\exp\!\left(\frac{\ln k}{\ell}t\right),\\[-1pt]
 n(t)&\approx n_0\exp\!\left(\frac{k-1}{\ell}t\right)
 \quad (k\approx1),\\[-1pt]
 \frac{dn}{dt}&\approx\frac{k-1}{\ell}n(t).
\end{aligned}
\]
Here $k$ is the generation-to-generation multiplication factor and $\ell$ is
the neutron lifetime used in this simple model.
\end{fixedbox}

\begin{boardbox}{Build the generation picture}
\centering
\begin{tikzpicture}[x=0.45in,y=0.27in,>=Latex,
  neutron/.style={circle,fill=KSUPurple,inner sep=1.5pt},
  fiss/.style={circle,draw=KSUGray,minimum size=0.22in}]
  \node[neutron] (n0) at (0,0) {};
  \node[fiss] (f0) at (0.7,0) {};
  \node[neutron] (n11) at (1.4,0.42) {};
  \node[neutron] (n12) at (1.4,-0.42) {};
  \node[fiss] (f11) at (2.1,0.42) {};
  \node[fiss] (f12) at (2.1,-0.42) {};
  \node[neutron] at (2.9,0.78) {};
  \node[neutron] at (2.9,0.30) {};
  \node[neutron] at (2.9,-0.30) {};
  \node[neutron] at (2.9,-0.78) {};
  \draw[->] (n0)--(f0);
  \draw[->] (f0)--(n11);
  \draw[->] (f0)--(n12);
  \draw[->] (n11)--(f11);
  \draw[->] (n12)--(f12);
  \draw[->] (f11)--(2.9,0.78);
  \draw[->] (f11)--(2.9,0.30);
  \draw[->] (f12)--(2.9,-0.30);
  \draw[->] (f12)--(2.9,-0.78);
  \node[below=0.10in of n0] {generation 0};
  \node[below=0.43in of n12] {generation 1};
  \node[below=0.10in] at (2.9,-0.78) {generation 2};
\end{tikzpicture}

\RaggedRight
Which neutrons count toward $k$? Why is ``neutrons emitted per fission'' not
itself the reactor multiplication factor?
\RuledLines{1}
\end{boardbox}


\begin{boardbox}{Criticality table}
\begin{tabularx}{\linewidth}{@{}c >{\bfseries}l X@{}}
\toprule
condition & name & neutron population\\
\midrule
$k<1$ & \RevealBlank[0.82in]{} & \RevealBlank[1.22in]{}\\[2pt]
$k=1$ & \RevealBlank[0.82in]{} & \RevealBlank[1.22in]{}\\[2pt]
$k>1$ & \RevealBlank[0.82in]{} & \RevealBlank[1.22in]{}\\
\bottomrule
\end{tabularx}
\end{boardbox}

\begin{checkpoint}[Connect neutron population to power]
\[
\begin{aligned}
 P(t)&=\MathBlank[0.52in]{}\,\MathBlank[0.70in]{},\\[-1pt]
 \frac{P(t)}{P(0)}&\approx\MathBlank[1.05in]{}.
\end{aligned}
\]
Why can even a small departure of $k$ from one matter when $\ell$ is short?
\RuledLines{1}
\end{checkpoint}

\begin{takeawaybox}[The three controls]
$E_f$ sets energy per fission; $k$ sets the generation-to-generation change;
$\ell$ sets the time scale. Together they set the scale and rate of energy
production.
\end{takeawaybox}

\end{handoutcolumns}
\end{document}
